% =====================================================================
\chapter[Radial Basis Function Neural Network]{Radial Basis Function\\Neural Network}
% =====================================================================

\labinfo{09\_RBF\_Network}{rbf\_network.py}{UCI Seeds: 210 wheat kernels, 7 geometric measurements,
3 varieties}{3-class classification}{Test accuracy 96.8\%, macro F1 0.968 (8 Gaussian hidden units)}

\section{Dataset}
The \textbf{Seeds} dataset (UCI) contains 210 wheat kernels, 70 from each of three varieties: Kama,
Rosa and Canadian. A soft X-ray image of each kernel was used to measure seven geometric properties:
area, perimeter, compactness ($4\pi\cdot\text{area}/\text{perimeter}^2$), kernel length, kernel width,
asymmetry coefficient and length of the kernel groove. All seven are continuous. I standardised them and
used a stratified 70/30 split (147 training and 63 test kernels).

\section{Code and experiments}
\begin{itemize}
\item \code{kmeans} is a plain implementation with 10 random restarts.
\item \code{RBFNetwork.fit} runs k-means (or, with \code{centre\_method="random"}, simply picks
      random training points), applies the width rule, builds the hidden-activation matrix with
      \code{\_hidden} and solves the least-squares problem for the output weights. A tiny ridge term
      ($10^{-8}$) keeps the solve stable when two hidden units are almost identical.
\item \code{predict} takes the arg-max of the outputs.
\end{itemize}
\noindent Experiments: (1)~a grid over the number of centres $K\in\{2,3,5,8,10,15,20,30,40\}$ and the
width scale $\in\{0.5,1,2,4\}$ with 5-fold CV; (2)~the best network on the test set; (3)~the same
network with random centres instead of k-means; (4)~a picture of a 2-D version of the network.

\section{Results}

\begin{table}[H]
\centering\small
\caption{5-fold cross-validation accuracy for different numbers of centres and widths.}
\label{tab:rbf-grid}
\begin{tabular}{@{}l c c c c c c c c c@{}}
\toprule
$K$ & 2 & 3 & 5 & 8 & 10 & 15 & 20 & 30 & 40\\
\midrule
width ×0.5 & 0.599 & 0.891 & 0.884 & 0.893 & 0.892 & 0.878 & 0.870 & 0.870 & 0.831\\
width ×1   & 0.769 & 0.925 & 0.925 & 0.905 & 0.905 & 0.905 & 0.891 & 0.899 & 0.904\\
width ×2   & 0.816 & 0.926 & 0.926 & 0.933 & 0.925 & 0.926 & 0.932 & 0.925 & 0.919\\
width ×4   & 0.831 & 0.926 & 0.919 & \textbf{0.953} & \textbf{0.953} & 0.919 & 0.945 & 0.952 & 0.939\\
\bottomrule
\end{tabular}
\end{table}

Two centres are clearly not enough for three classes. From three centres on, accuracy is above 89\% for
every width. Narrow Gaussians (×0.5) are consistently worst: each unit only ``sees'' a small region and
the network starts to memorise individual kernels. The widest setting worked best, which tells us that
the three varieties are separated by smooth, gently curved boundaries. The best result was $K=8$ with
width scale 4 ($\sigma=6.72$ in standardised units), with a CV accuracy of 95.3\%.

\begin{table}[H]
\centering\small
\caption{Test-set results of the final network ($K=8$, $\sigma=6.72$, 63 kernels).}
\label{tab:rbf-test}
\begin{tabular}{@{}l S[table-format=1.3] S[table-format=1.3] S[table-format=1.3] S[table-format=2.0]@{}}
\toprule
\textbf{Variety} & {\textbf{Precision}} & {\textbf{Recall}} & {\textbf{F1}} & {\textbf{Support}}\\
\midrule
Kama     & 0.952 & 0.952 & 0.952 & 21\\
Rosa     & 1.000 & 1.000 & 1.000 & 21\\
Canadian & 0.952 & 0.952 & 0.952 & 21\\
\midrule
Macro average & 0.968 & 0.968 & 0.968 & 63\\
\midrule
\multicolumn{5}{@{}l}{Accuracy: \textbf{0.968} (61 of 63)}\\
\multicolumn{5}{@{}l}{Random centres: 0.968 (width ×4), 0.873 (width ×1, where k-means gives 0.905)}\\
\bottomrule
\end{tabular}
\end{table}

The network classified 61 of the 63 test kernels correctly. Kama sits between the other two varieties
in size, so its border with Canadian is where the mistakes happen. With these wide Gaussians, the
exact positions of the centres hardly mattered: random centres gave the same test accuracy as k-means
(96.8\%). With the plain width rule (scale ×1) the centres mattered more: k-means reached 90.5\% and
random centres only 87.3\%, because narrow bumps placed in the wrong spots leave parts of the input
space uncovered.

\begin{figure}[H]
\centering
\includegraphics[width=0.48\linewidth]{\figpath{09\_RBF\_Network}{cv_vs_centres.pdf}}\hfill
\includegraphics[width=0.48\linewidth]{\figpath{09\_RBF\_Network}{rbf_2d_regions.pdf}}
\caption{RBF centre/width selection and the two-dimensional decision regions.}
\end{figure}

\section{Source code}
\lstinputlisting[caption={\texttt{09\_RBF\_Network/rbf\_network.py}}]{09_RBF_Network/rbf_network.py}
\clearpage
